\subsection*{Grounding Chinese reply images in a printed catalogue}
We used the official dialogue splits and venue databases of \href{https://github.com/thu-coai/CrossWOZ}{CrossWOZ}. We retained single-field human requests with a matching informed value, a resolved venue and at least one competing attribute. We required the exact catalogue value in the human answer and bounded input geometry. This yielded 5,097 of 42,346 training system turns, 517 validation turns and 515 of 4,238 test system turns. The model received a pictured natural question and a supplied venue record in a known typeface and fixed 32-pixel cell layout. We supplied the resolved venue to both readers. Whole questions and venues may recur across the official dialogue splits.

Field prefixes delimit contiguous source pixels. Given the query image, domain and available field mask, the encoder predicts the requested field. Training supervises field identity and the likelihood of the human answer. A causal writer chooses among 437 printable training glyphs, a name-copy action, a selected-value-copy action and stop. Copy actions insert the original cell sequence. Optional transcription follows image generation. We matched the byte-input control on the 8,000 updates, 2,998,698 target actions and physical answer task. Each network has 1,568,944 parameters. Device energy for fitting was 10.57 and 9.91 kJ for image and byte input, respectively. These timings exclude corpus preparation. Malformed phrasing and unsupported qualitative claims can pass the value test.

For a separate scope test, we held card content fixed and exchanged two field names in the question. A bidirectional recurrent selector read the glyph context at two equalized field markers and chose one of their positions. We used 20 labelled sentence frames with 24 prefix/suffix combinations and trained for 1,000 steps (2.98 seconds, 0.27 kJ). Six field names occurred in fitting, with two held out. We tested the fixed original frames and twelve further constructions, each on 18 field pairs. The 216 pairs share twelve linguistic frames. Exact mention detection activates this selector only for two distinct visible labels.
